	\chapter{Altri approcci alla computabilità e tesi di Church}
	
	\section{Altri approcci alla computabilità}
	
	Il modello URM è uno dei modelli più recenti usati nello studio della calcolabilità. Tra gli approcci alla calcolabilità, questi sono alcuni dei modelli più noti:
	\begin{itemize}
		\item Godel - Herbrand - Kleene: funzioni ricorsive generali.
		\item Church: $\lambda$-calcolo.
		\item Godel - Kleene: funzioni $\mu$-ricorsive e funzioni ricorsive parziali. $\star$
		\item Turing: funzioni calcolabili mediante macchine di Turing. $\star$
		\item Post: funzioni definite mediante sistemi deduttivi canonici.
		\item Sheperdson-Sturgis: funzioni URM-calcolabili. $\star$
	\end{itemize}
	
	i modelli $\star$ saranno oggetti del nostro studio.
	
	\paragraph{Risultato della teoria della calcolabilità} Tutti i modelli precedentemente proposti caratterizzano la medesima classe di \textbf{funzioni calcolabili}, che abbiamo precedentemente denotato con $\mathcal{C}$, per cui non ha più senso dire $\mathcal{C}_\text{URM}$.
	
	\section{Funzioni ricorsive parziali (Godel - Kleene)}
	La classe $\mathcal{R}$ delle \textbf{funzioni ricorsive parziali} è la più piccola delle famiglie di funzioni parziali che
	\begin{itemize}
		\item Contiene le funzioni iniziali $\underline 0, x+1, \bigcup_{i}^n$.
		\item È chiusa rispetto a \textbf{composizione}, \textbf{ricorsione primitiva}, \textbf{minimalizzazione}.
	\end{itemize}
	
	Godel e Kleene hanno inizialmente portato la loro attenzione sulla classe $\mathcal{R}_0$, ossia delle \textbf{funzioni $\mu$-ricorsive}, che coincide con $\mathcal{R}$ e ammette la minimalizzazione \textbf{solo se non consente funzioni parziali}.
	
	$$
	\mathcal{R}_0 = \mathcal{R} \cap TOT
	$$
	
	\subsection{Tutte le funzioni ricorsive parziali sono calcolabili}
	$$
	\mathcal{R} = \mathcal{C}_{URM}
	$$
	
	\paragraph{Dimostrazione} 
	
	Per quanto visto in precedenza, $\mathcal{R} \subseteq \mathcal{C}_{URM}$.

	Viceversa, sia $f(\vec{x})$ una funzione URM-calcolabile da un programma $P = I_1, I_2, \dots, I_s$. Senza perdere di generalità, possiamo supporre $P$ sia in forma standard, e che non ci siano istruzioni di salto che facciano passare all'istruzione subito successiva. Definiamo quindi un \textit{passo} come l'esecuzione di un'istruzione. 
	
	Consideriamo adesso le seguenti le seguenti funzioni relative alla computazione di $P$.
	
	$$
	c(\vec{x}, t) = \begin{cases}
	R_1 &\text{dopo $t$ passi della computazione $P(\vec{x})$ se questa non si è interrotta.}\\
	R_1 &\text{a fine computazione, se si è interrotta.}
	\end{cases}
	$$
	
	$$
	j(\vec{x}, t) =  \begin{cases}
		\text{numero della prossima istruzione} &\text{dopo $t$ passi di $P(\vec{x})$ se questa non si è interrotta.}\\
		0 &\text{se si è interrotta.}
	\end{cases}
	$$
	
	queste sono evidentemente funzioni totali. Osserviamo quindi che
	
	$$
	f_P^{(n)}(\vec{x})\downarrow \Leftrightarrow \mu t (j(\vec))
	$$
NON STO CAPENDO NULLA
	
	\section{Macchine di Turing}
	Una macchina di Turing $M$è un dispositivo finito che effettua operazioni primitive su un nastro infinito. Il nastro infinito è suddiviso in caselle che possono contenere un simbolo tratto da una lista finita di simboli, detto \textbf{alfabeto della macchina}. $s_0$ è il simbolo \textit{"blank"}, o $B$.
	
	$M$ può effettuare operazioni di lettura e scrittura (tramite una testina), e spostarsi a sinistra o a destra sul nastro. Ha inoltre un'unità di controllo costituita da un insieme finito di quadruple.
	
	\paragraph{Quadruple} Ciascuna di queste quadruple è nella forma
	
	$$
	q_i \ s_j \ \alpha \ q_l
	$$
	
	con $q_i$ stato della macchina, $s_j$ carattere letto, $\alpha \in \{L, R, s_0, \dots, s_k\}$ e $q_l$ nuovo stato. Ogni quadrupla definisce quindi quale azione effettuare (spostamento a sx, a dx o sovrascrittura di quale simbolo), quando la macchina legge un carattere in un determinato stato, e in quale stato transitare.
	
	\paragraph{Convezione sul contare} Negli usi che faremo, rappresenteremo il numero $n$ con $n$ caselle contenenti il simbolo $1$. La funzione che conta il numero di occorrenze di $1$ sul nastro è Turing-calcolabile.
	
	\subsection{Funzioni Turing calcolabili}
	Una funzione parziale si dice \textbf{Turing-calcolabile} se esiste una macchina di Turing che la calcola. La famiglia delle funzioni Turing-calcolabili si indica con $\mathcal{C}_{TUR}$.
	
	\subsection{Equivalenza tra URM e TM}
	Le funzioni ricorsive parziali $\mathcal R$ sono URM-calcolabili, ma anche Turing-calcolabili, ossia
	
	$$
	\mathcal{C}_{URM} = \mathcal{R}= \mathcal{C}_{TUR}
	$$
	
	e ciò è dimostrabile per \textbf{simulazione}.
	
	\newpage
	
	\section{Tesi di Church-Turing}
	La tesi di Church Turing è uno dei pilastri teorici dell'informatica moderna, e ha un ruolo fondamentale nella teoria della computabilità.
	Afferma che:
	\begin{center}
	\textit{La classe delle funzioni \textbf{intuitivamente} (umanamente) calcolabili coincide con la classe $\mathcal{C}_{URM}$, ossia quella delle funzioni URM-calcolabili. Lo stesso vale per le macchine di Turing, e qualsiasi altro sistema formale di computazione. È calcolabile tutto ciò che può essere calcolato da una macchina di Turing.}
	\end{center}
	
	A supporto di questa tesi, abbiamo che:
	\begin{itemize}
		\item La tesi di Church-Turing risulta essere il teorema fondamentale della calcolabilità, in quanto tutti i sistemi di calcolabilità sembrano convergere alla stessa classe di funzioni calcolabili.
		\item Si ha una collezione di funzioni per la quale esiste una dimostrazione di appartenenza a $\mathcal{C}_URM = \mathcal{C}_TUR = \mathcal{R}$.
		\item La possibilità di simulare effettivamente programmi URM.
		\item La mancanza di controesempi di funzioni effettivamente\footnote{Nella teoria della calcolabilità, un algoritmo effettivo (o procedura effettiva) è un insieme finito di istruzioni precise, non ambigue e meccaniche che, se applicato a un problema, produce una soluzione corretta in un numero finito di passi. Deve essere eseguibile da un essere umano o macchina senza intuizione, garantendo lo stesso risultato.} calcolabili ma non URM-calcolabili.
	\end{itemize}
	
	\subsection{La potenza della tesi di Church-Turing}
	Basandosi sulle prove precedentemente descritte, tutti i matematici hanno accettato la tesi di Church, la quale risulta essere uno strumento estremamente potente nelle dimostrazioni. Piuttosto che usare un modello teorico, come quello URM, o le macchine di Turing, possiamo procedere nel seguente modo.
	
	\paragraph{Dimostrazione della calcolabilità \textit{per tesi di Church}} Desideriamo dimostrare la calcolabilità della funzione $f$. Definiamo (in linguaggio naturale, tramite diagrammi o schemi) un algoritmo, e forniamo una dimostrazione informale ma rigorosa relativa al fatto che l'algoritmo calcola effettivamente la funzione $f$. Appelliamoci poi alla tesi di Church e concludiamo immediatamente dicendo che $f$ è URM-calcolabile. 
	
	\subsubsection{Estratto dal libro (che personalmente adoro) riguardo la tesi di Church-Turing}

	\begin{verbatim}
	Note that in using Church’s thesis we are not proposing to abandon all
	thought of proof, as if Church’s thesis is a magic wand which we can wave instead. 
	A proof by Church’s thesis will always involve proof that is
	careful, and sometimes complicated, although informal. [...]
	Church’s thesis not only keeps proofs shorter, but also prevents the
	main idea of a proof or construction from being obscured by a mass of
	technical details. It remains, however, an expression of faith or
	confidence. The validity of faith depends on the evidence that can be
	mustered. In the case of Church’s thesis, there is the mathematical
	evidence already outlined. For the practised student there is the addi
	tional evidence of his own experience in translating informal algorithms
	into formal counterparts. For the beginner, our use of Church’s thesis in
	subsequent chapters may call on his willingness to place confidence in the
	ability of others until self confidence is developed.
	\end{verbatim}
	